\documentclass{article}
\usepackage[utf8]{inputenc}
\usepackage[T1]{fontenc}
\usepackage{geometry}
\geometry{
    a4paper,
    left=2.5cm,
    right=2.5cm,
    top=2.5cm,
    bottom=2.5cm
}
\usepackage{graphicx}
\usepackage{float}
\usepackage{hyperref}
\usepackage[english]{babel}
\usepackage{amsmath}
\usepackage{amssymb}
\usepackage{array}
\usepackage{booktabs}
\usepackage{caption}
\usepackage{subcaption}
\usepackage{enumitem}
\usepackage{fancyhdr}
\usepackage{titlesec}
\usepackage{url}
\usepackage{textcomp}
\usepackage{xcolor}
\usepackage{listings}
\usepackage{multirow}
\usepackage{longtable}
\usepackage{tabularx}
\usepackage{times}

% Set up page style
\pagestyle{fancy}
\fancyhf{}
\rfoot{\thepage}
\renewcommand{\headrulewidth}{0pt}

% Title formatting
\titleformat{\chapter}[display]{\normalfont\bfseries}{\chaptertitle\thechapter}{0pt}{\Huge}
\titlespacing*{\chapter}{0pt}{50pt}{40pt}

% Document info
\title{Template Document}
\author{Academic Support Agent}
\date{\today}

\begin{document}
\maketitle

% ----------------  TITLE PAGE ------------------- %
% no need of \maketitle, as title-page is added
\noindent 
% Content from Sections/titlepage would be included here
 %for Image to work, changed the compiler to pdfLatex
\newpage

% ---------------- CONTENT PAGES ----------------- %
\pagenumbering{Roman}
% TOC

% Use "\addcontentsline{toc}{section}{name}" before name to put in toc
\newpage

% List of tables
\addcontentsline{toc}{section}{List of Tables} %to add to table of contents

% If not automatically added -- \addcontentsline{lot}{section}{Table 1}
\newpage

% List of Figures
\addcontentsline{toc}{section}{List of Figures}

% If not automatically added -- \addcontentsline{lof}{section}{Fig 1}
\newpage

% List of acronyms 
\addcontentsline{toc}{section}{Acronyms}
\printglossary[type=\acronymtype]

% List of Symbols 
\addcontentsline{toc}{section}{List of Symbols}
\printglossary[type=symbols,style=long,title={List of Symbols}]

\newpage

% ------------------  ABSTRACT ---------------------- %
\pagenumbering{arabic}
\addcontentsline{toc}{section}{Abstract}
\section*{Abstract} 
This report aims to evaluate xx. The abstract should be double-spaced and not more than 100 words

% -------------------------------- MAIN PART ----------------------------- %

% ------------------ INTRODUCTION ------------------- %
% Content from Sections/Introduction would be included here


% --------------------  BODY  ----------------------- %
% Content from Sections/Body would be included here


% -----------------  CONCLUSION  --------------------- %
% Content from Sections/Conclusion would be included here



\section{ALL SAMPLE DATA AND EGS}
\subsection{subsec 0}
This is some text for the section.

I can reference now otherwise use (\citet{dirac}). Now trying the footnotes. 
It is pretty simple \footnote{Some UTF-8 support via \LaTeX{} kernel commands.}

\subsection{subsec 1}
An analysis of the effect of the revised recruitment policies on staff turnover at divisional headquarters (\citet{einstein})

The equation given by einstein - 
\begin{equation}\label{Emc2}
E=mc^2
\end{equation}

Einstein wrote his famous equation \eqref{Emc2}, blah blah blah ....

Integrals
\begin{align*}
  f(x) &= x^2\\
  g(x) &= \frac{1}{x}\\
  F(x) &= \int^a_b \frac{1}{3}x^3
\end{align*}

\subsection{This is an example of subsection Title}
Reference symbols: $\gls{x}$, $\gls{v}$, $\gls{a}$, $\gls{t}$, $\gls{F}$.
Given a set of numbers, there are elementary methods to compute its \glsxtrlong{gcd}, which is abbreviated \glsxtrshort{gcd}. This process is similar to that used for the \glsxtrfull{lcm}.

Using another footnote here \footnote{some data}

\subsubsection{subsubsec1:(\texttt{ragged2e})}

\begin{Center}
\blindtext[1]
\end{Center}

\subsubsection{subsubsec 2}
\Centering
\blindtext[1]

\vspace{5mm}
\justifying

\blindtext[1]
\lipsum[2]

% TO SET PARAGRAPH INDENTING AND SPACING
\setlength{\parindent}{15pt} % for indenting the paras
\setlength{\parskip}{0.5em} % for spacing between paras % em being a different unit

\subsection*{subsec 3}
\textbf{First paragraph} of a section which, as you can see, is not indented. This is more text in the paragraph. This is more text in the paragraph.

\textbf{Second paragraph}. As you can see it is indented. This is more text in the paragraph. This is more text in the paragraph. 

\noindent\textbf{Third paragraph}. This too is not indented due to use of \texttt{\string\noindent}. This is more text in the paragraph. This is more text in the paragraph.  The current value of \verb|\parindent\raggedright| is \the\parindent. This is more text in the paragraph.

\subsection*{subsec 4}
\textbf{First paragraph} of a section which, as you can see, is not indented. This is more text in the paragraph. This is more text in the paragraph.
\lipsum[]

\subsection{subsec 5}

Tables:

In text, I can reference the Table \ref{tab:my_label1} like this.

%\setcounter{table}{0} % sets table counter to zero 
\begin{table}[ht]
    \centering
    \begin{tabular}{|c||c|c|c|}
    \hline
    $x$     & 1  & 2  & 3  \\ 
    \hline
    $f(x)$  & 10 & 12 & 14 \\ 
    \hline
    \end{tabular}
    \caption{e.g. Table 1}
    \label{tab:my_label1}
\end{table}

Rows and columns can be merged to create larger table cells. 

\begin{table}[ht]
\begin{center}
\begin{tabular}{ |c|c|c|c| } 
\hline
col1 & col2 & col3 \\
\hline
\multirow{3}{4em}{Multiple row} & cell2 & cell3 \\ 
& cell5 & cell6 \\ 
& cell8 & cell9 \\ 
\hline
\end{tabular}
\end{center}
\caption{e.g. Table 2}
\label{tab:my_label2}
\end{table}

The table \ref{table:3} and \ref{table:4} are an examples of referenced \LaTeX{} elements.

\begin{table}[ht]
\centering
\begin{tabular}{||c c c c||} 
 \hline
 Col1 & Col2 & Col2 & Col3 \\ [0.5ex] 
 \hline\hline
 1 & 6 & 87837 & 787 \\ 
 2 & 7 & 78 & 5415 \\
 3 & 545 & 778 & 7507 \\
 4 & 545 & 18744 & 7560 \\
 5 & 88 & 788 & 6344 \\ [1ex] 
 \hline
\end{tabular}
\caption{This is the caption for the third table.}
\label{table:3}
\end{table}

\begin{table}[ht]
\centering
\begin{tabular}{||c c c c||} 
 \hline
 Col1 & Col2 & Col2 & Col3 \\ [0.5ex] 
 \hline\hline
  4 & 545 & 18744 & 7560 \\
 5 & 88 & 788 & 6344 \\ [1ex] 
 \hline
\end{tabular}
\caption{This is the caption for the fourth table.}
\label{table:4}
\end{table}

Images
\begin{figure}[ht]
\centering
    \begin{subfigure}{0.45\textwidth}
        \includegraphics[width=0.9\linewidth, height=6cm]{overleaf-logo} 
        \caption{img 1}
        \label{fig:subim1}
    \end{subfigure}
    \begin{subfigure}{0.5\textwidth}
        \includegraphics[width=0.9\linewidth, height=6cm]{mesh}
        \caption{img 2}
        \label{fig:subim2}
    \end{subfigure}
\caption{Figure with two images}
\label{fig:fig1}
\end{figure}

I can reference images such as \ref{fig:fig1} containing \ref{fig:subim1} and \ref{fig:subim2}

%Floatting text around a figure
Praesent in sapien. Lorem ipsum dolor sit amet, consectetuer adipiscing elit. Duis fringilla tristique neque. Sed interdum libero ut metus. Pellentesque placerat.

\begin{wrapfigure}{l}{0.25\textwidth}
\includegraphics[width=0.9\linewidth]{overleaf-logo} 
\caption{img3}
\label{fig:fig2}
\end{wrapfigure}

Praesent in sapien. Lorem ipsum dolor sit amet, consectetuer adipiscing elit. Duis fringilla tristique neque. Sed interdum libero ut metus. Pellentesque placerat. Nam rutrum augue a leo. Morbi sed elit sit amet ante lobortis sollicitudin. (See figure \ref{fig:fig2} )

Praesent in sapien. Lorem ipsum dolor sit amet, consectetuer adipiscing elit. Duis fringilla tristique neque. Sed interdum libero ut metus. Pellentesque placerat. Nam rutrum augue a leo. Morbi sed elit sit amet ante lobortis sollicitudin.Praesent in sapien. Lorem ipsum dolor sit amet, consectetuer adipiscing elit. Duis fringilla tristique neque. Sed interdum libero ut metus. Pellentesque placerat. Nam rutrum augue a leo. Morbi sed elit sit amet ante lobortis sollicitudin. Pellentesque placerat. Nam rutrum augue a leo. Morbi sed elit sit amet ante lobortis sollicitudin.

% ------------ APPENDIX ------------------ %



\section{App / Anx 1 - Formula XXX}
Text and images for Ap A.
The appendix is an optional part of the thesis, which might include: tables with additional data, additional figures, mathematical proofs and derivations, etc. The different parts of the appendix should be lettered (e.g. Appendix A") and might be subdivided (e.g. A.1"). The appendices are listed in the table of contents
(see example documents). You have to refer to each appendix at least once in the main part, otherwise this appendix serves no purpose.

\section{App / Anx 2 - Report XXX}
Report data and lab results or prog methods for Ap B

% ------------ Bibliography ------------------ %
\addcontentsline{toc}{section}{Bibliography and references}


% ------------ Author Declaration ------------ %
\addcontentsline{toc}{section}{Declaration of Author}
% Content from Sections/AuthorsDeclaration would be included here
\end{document}